\section{A life written through books}\label{sec:aug-writings}

\subsection{From dialogues to pastoral interpretation}

Augustine's books are events in his biography. The Cassiciacum dialogues of
386--387 display a convert testing how classical disciplines, philosophical
argument, and Christian authority might coexist. \emph{On the Immortality of
the Soul}, \emph{On Music}, and \emph{On the Teacher} continue questions about
mind, signs, learning, and inward illumination. They do not yet contain every
later formulation, and their dramatic speakers cannot always be equated
without remainder with Augustine's settled voice.

After ordination, Scripture and controversy increasingly governed the work.
\emph{On Christian Teaching}, composed in more than one stage,
joined interpretation with formation in charity and instruction in effective
communication. The famous distinction between things and signs is not an
invitation to decode texts apart from the Church's rule of faith and the
double command of love. Augustine also authorized the use of learning found
outside Christianity, comparing it to Israel's appropriation of Egyptian
goods (\emph{On Christian Teaching} 2.40.60). The image has generated a long
tradition of Christian humanism, though it belongs to a late-antique argument,
not a modern theory of neutral disciplines.

His expositions of the Psalms and tractates on John grew from preaching and
revision. They show a characteristic habit: Christ speaks as head and body,
so the psalm can become the prayer of the whole Christ. This ecclesial reading
connects exegesis, sacrament, poverty, martyrdom, and moral transformation.
Augustine's interpretive imagination is often more visible here than in
isolated doctrinal propositions.

\subsection{The \emph{Confessions}: memory addressed to God}

The thirteen books of the \emph{Confessions}, composed after Augustine became
bishop, are not simply the first nine books of autobiography. Books 1--9
remember the journey to baptism and Monnica's death. Book 10 examines memory,
temptation, and the still-unfinished converted self. Books 11--13 interpret
Genesis, creation, time, and the Church. The movement prevents the reader from
treating conversion as the end of moral struggle. The same “I” who narrates
the garden confesses present distraction and asks how a temporal creature
hears the eternal Word.

The work's truth is inseparable from its addressee. Augustine speaks to God
before the reader. Praise, confession of sin, confession of faith, biblical
allusion, remembered dialogue, and philosophical inquiry interpenetrate.
Modern biographers may recover much from it, but a seamless factual digest
destroys the form of the evidence. The pear theft, Monnica's visions, the
voice in the garden, and the Ostia conversation are remembered because they
serve an argument about providence and desire.

\subsection{\emph{The Trinity}, grace, and inward relation}

Augustine worked on \emph{The Trinity} across many years; this revision does
not assign the individual books an independently checked modern chronology.
Early books defend Nicene faith and interpret
Scripture; later books seek images of triune relation in the human mind.
Memory, understanding, and will are not a proof that contains God. They are
created vestiges whose inadequacy becomes part of the ascent. The work's long
composition matters: it belongs to a bishop interrupted by controversy,
administration, and unauthorized circulation, not to an isolated speculative
career.

His theology of grace likewise cannot be detached from prayer. In
\emph{Confessions} 10.29.40 he asks God to grant what God commands and then to
command what God wills. Its logic is already biographical: the converted
person obeys genuinely, but obedience itself is enabled by gift. Later
anti-Pelagian works sharpened the
priority of grace and the mystery of election. Catholic reception drew deeply
from this teaching while continuing to debate how best to speak of freedom,
predestination, and God's universal saving will.

\subsection{\emph{The City of God} after the sack of Rome}

Alaric's sack of Rome in 410 shocked the Mediterranean and sent refugees to
Africa. Some blamed Christianity for weakening Rome's ancestral protections.
Augustine began \emph{The City of God} after the sack and completed its
twenty-two-book argument before reviewing the work in \emph{Retractationes}
2.43; a finer year-by-year chronology was not controlled here. The books exceed
the occasion. The first part criticizes pagan
claims about divine protection and philosophical beatitude; the second traces
two cities formed by different loves through biblical history toward distinct
ends.

The “earthly city” is not simply the state, and the “city of God” is not
identical in every respect with the visible Church. The two are intermingled
in history. Political peace is a real, limited good that Christians can share;
no empire becomes the eschatological kingdom by success. Yet Augustine did not
write a blueprint for liberal democracy, theocracy, or cultural withdrawal.
Later political Augustinianisms select and extend different elements. A
biography should locate the book within refugee crisis, episcopal government,
Roman fragility, and Augustine's biblical horizon before recruiting it into
later ideologies.

\subsection{Speech, correspondence, and collaborative production}

Augustine's surviving corpus depended on stenographers, scribes, messengers,
libraries, patrons, readers, and clerical colleagues. Letters could become
public dossiers. Sermons might circulate without final authorization. Books
were dictated, revised, excerpted, and transported across dangerous roads and
sea lanes. Possidius's remark that Augustine left the Church books and a
library (\emph{Life} 31) gives material weight to the intellectual legacy.

His correspondence reveals a range hidden by the canonical masterpieces:
questions from lay readers, requests for intercession, disciplinary crises,
arguments over biblical translation, and anxious efforts to keep communities
together. Modern editions include collections commonly called the Divjak
letters and Dolbeau sermons, but their contents were not collated for this
biography and are not used to support a new portrait. The already cited
correspondence is sufficient reason to resist a hierarchy in which the “real”
Augustine is only the author of a few timeless books and pastoral paperwork is
background noise.

\subsection{Author, authority, and Doctor of the Church}

Augustine's authority grew during his life through networks and controversy.
It grew after death through manuscript copying, florilegia, monastic rules,
theological disputes, councils, liturgy, medieval schools, Reformation
controversy, modern philosophy, and papal teaching. The Catholic title Doctor
of the Church receives his eminent service to doctrine; it does not canonize
each historical inference or rhetorical strategy. Benedict XVI's audience of
9 January 2008 presents Augustine as a companion in the search for truth and
as a theologian whose conversion remained lifelong. That is official reception
(M), not independent evidence for the precise date of a fourth-century event.

Augustine's own practice offers a more demanding model of authority than
quotation collecting. He argued, changed, re-read, and submitted his work to
the truth he believed exceeded him. The most faithful historical reading is
therefore neither automatic agreement nor dismissal. It identifies which
Augustine is speaking, to whom, under what pressure, and with what later
correction.
